% ============================================================
%  LEMBAR HASIL REVISI SIDANG SKRIPSI
%  Dokumen mandiri (standalone) — bukan bagian dari thesis.tex.
%  Compile:  pdflatex lembar-hasil-revisi-sidang.tex
%
%  Kolom "Catatan Revisi Penguji" diisi verbatim dari catatan penguji di
%  SISIDANG (REVISI BERKAS, Revisi #1), didistilasi pada issue #76-#84
%  berlabel `revisi-sidang`. Kolom "Hasil/Jawaban Revisi" dan "Halaman Hasil
%  Revisi" sengaja dikosongkan untuk diisi tangan setelah revisi rampung.
% ============================================================
\documentclass[12pt, a4paper]{article}
\usepackage{iftex}

% -- Encoding, Language & Fonts --------------------------------
\ifXeTeX
    \usepackage{polyglossia}
    \setmainlanguage{bahasai}
    \usepackage{fontspec}
    \setmainfont{TeX Gyre Termes}
    \setsansfont{TeX Gyre Heros}
\else
    \usepackage[utf8]{inputenc}
    \usepackage[T1]{fontenc}
    \usepackage[indonesian]{babel}
    \usepackage{mathptmx}
    \usepackage{helvet}
\fi

% -- Page geometry ---------------------------------------------
\usepackage[a4paper, top=2.5cm, bottom=2.5cm, left=3cm, right=2.5cm]{geometry}

% -- Spacing & paragraph ---------------------------------------
\usepackage{setspace}
\onehalfspacing
\setlength{\parindent}{0pt}
\setlength{\parskip}{4pt}

% -- Tables ----------------------------------------------------
\usepackage{longtable, array}
\usepackage{microtype}

% -- Header/footer ---------------------------------------------
\usepackage{fancyhdr}
\pagestyle{fancy}
\fancyhf{}
\fancyfoot[C]{\small\thepage}
\renewcommand{\headrulewidth}{0pt}

\usepackage[hidelinks]{hyperref}

% ============================================================
\begin{document}

\begin{center}
  {\large\bfseries LEMBAR HASIL REVISI SIDANG SKRIPSI}\\[2pt]
  {\normalsize Program Studi Sarjana Ilmu Komputer\\
  Fakultas Ilmu Komputer, Universitas Indonesia}
\end{center}
\vspace{0.4cm}

% -- Identitas ------------------------------------------------
\noindent\begin{tabular}{@{}p{3.2cm}@{ : }p{\dimexpr\linewidth-3.8cm}@{}}
Judul Skripsi &
LLM-Assisted Knowledge Graph Completion untuk Pemetaan Interkoneksi Materi
Sains (Fisika, Kimia, Biologi) pada Kurikulum Merdeka SMA Kelas XII \\[6pt]
Mahasiswa &
Fadrian Yhoga Pratama (2206819395); Soros Febriano (2206083445);
Tegar Wahyu Khisbulloh (2206082032) \\[6pt]
Program Studi & Sarjana Ilmu Komputer \\[6pt]
Pembimbing & Dr.\ Siti Aminah, S.Kom., M.Kom. \\[6pt]
Tanggal Sidang & \rule{5cm}{0.4pt} \\
\end{tabular}

\vspace{0.5cm}

% -- Tabel revisi ---------------------------------------------
\renewcommand{\arraystretch}{1.3}
\setlength{\tabcolsep}{5pt}
\small
\begin{longtable}{|c|p{6.4cm}|p{4.6cm}|c|}
\hline
\textbf{No} &
\multicolumn{1}{c|}{\textbf{Catatan Revisi Penguji}} &
\multicolumn{1}{c|}{\textbf{Hasil/Jawaban Revisi}} &
\multicolumn{1}{c|}{\parbox[c]{1.9cm}{\centering\textbf{Halaman Hasil Revisi}}} \\
\hline
\endfirsthead
\hline
\textbf{No} &
\multicolumn{1}{c|}{\textbf{Catatan Revisi Penguji}} &
\multicolumn{1}{c|}{\textbf{Hasil/Jawaban Revisi}} &
\multicolumn{1}{c|}{\parbox[c]{1.9cm}{\centering\textbf{Halaman Hasil Revisi}}} \\
\hline
\endhead
\hline
\multicolumn{4}{|r|}{\small\textit{Bersambung ke halaman berikutnya}} \\
\hline
\endfoot
\hline
\endlastfoot

% ===== Penguji 1 =====
\multicolumn{4}{|l|}{\textbf{Penguji 1 — Evi Yulianti, M.Comp.Sc., M.Kom., Ph.D.}} \\
\hline
1 & Rumusan masalah (RQ) perlu diperbaiki; jangan menanyakan cara. Cara ada di
Bab~3, sedangkan jawaban RQ ada di kesimpulan. & & \\
\hline
2 & Jelaskan \textit{task} apa yang dikerjakan; buat ilustrasi \textit{task}
yang menggambarkan input--output. & & \\
\hline
3 & Kenapa sistem ini dibuat? Kebutuhannya dari mana? Motivasi pentingnya
sistem perlu digali lagi. & & \\
\hline
4 & Pemanfaatan sistem ini untuk apa nanti? Bagaimana skenario
penggunaannya? & & \\
\hline
5 & Kenapa mata pelajaran Fisika, Biologi, dan Kimia yang dipilih? & & \\
\hline
6 & Siapa pakarnya? Apa saja yang divalidasi oleh pakar? & & \\
\hline
7 & Metode untuk menggabungkan \textit{triple} (\textit{merge} KG) harus
dijelaskan. & & \\
\hline
8 & Untuk \textit{triple} yang dinilai ``salah'' oleh pakar, apakah digunakan
lagi pada tahap~2 (\textit{merge} KG)? & & \\
\hline

% ===== Penguji 2 =====
\multicolumn{4}{|l|}{\textbf{Penguji 2 — Iis Afriyanti, S.Kom., M.Sc.}} \\
\hline
9 & Perbaiki rumusan masalah dan kaitkan dengan jawabannya di kesimpulan. & & \\
\hline
10 & Kenapa pengembangan ontologi tidak dijadikan sebagai pertanyaan
penelitian? & & \\
\hline
11 & Relasi yang salah atau tidak valid kenapa dibuang atau diikutkan untuk
tahap selanjutnya? Berikan argumentasinya. & & \\
\hline
12 & Perbaiki bagaimana cara mengembangkan ontologi. & & \\
\hline
13 & Tambahkan argumentasi kenapa memakai OWL? Kenapa bukan yang lain? Apa
yang dimanfaatkan dari OWL? & & \\
\hline
14 & Tambahkan argumentasi kenapa memilih Gemini Pro sebagai
\textit{LLM-as-a-judge}? & & \\
\hline

\end{longtable}

\vspace{0.3cm}
\noindent\small\textit{Catatan: kolom ``Catatan Revisi Penguji'' dikutip dari
catatan penguji pada SISIDANG (Revisi Berkas~\#1). Kolom ``Hasil/Jawaban
Revisi'' dan ``Halaman Hasil Revisi'' diisi setelah perbaikan naskah skripsi
selesai dilakukan.}

% -- Tanda tangan penguji -------------------------------------
\vspace{1.2cm}
\normalsize
\noindent\begin{tabular}{@{}p{0.5\linewidth}@{}p{0.5\linewidth}@{}}
\centering Penguji 1, & \centering Penguji 2, \tabularnewline
\\[2cm]
\centering\rule{6cm}{0.4pt} & \centering\rule{6cm}{0.4pt} \tabularnewline
\centering Evi Yulianti, M.Comp.Sc., M.Kom., Ph.D. &
\centering Iis Afriyanti, S.Kom., M.Sc. \tabularnewline
\end{tabular}

\end{document}
